\part{The research process}

\chapter{Experiments, like a research team}
\chapfig{fig_research_team.png}

A single person tuning one model at a time is slow and biased. We instead ran experiments the way a research
lab does: \textbf{many models in parallel, each cheap, the losers killed early, the winner promoted.}

\fig{figures/diag_experiments.png}{0.98}{The selection funnel. Seven ablations choose the learning rate and
mix; four full models then race on the data mix; the winner is promoted and finally scaled to 1.3B. Abandoned
experiments (grey) die early so they never waste a GPU-day.}

\section{Round one: the ablation sweep}
\fig{plots/plot_ablation.pdf}{0.82}{Seven short models trained at once (one GPU each, $\sim$0.35\,B tokens),
ranked by PubMed validation loss at a matched step. Learning rate \texttt{1e-3} edged out the field;
\texttt{3e-4} was clearly the laggard; grouped-query attention was essentially free.}

\noindent Seven models, varying learning rate, data mix, and one architecture choice (grouped-query
attention), all running simultaneously, each self-abandoning if it stalled. The finding: a learning rate of
\textbf{1e-3} --- higher than the textbook default --- and that GQA costs nothing, so we could keep it for
faster inference. Total cost of the round: about \textbf{\$15}.

\section{Round two: the production fleet}
\fig{plots/plot_fleet.pdf}{0.82}{Four full models trained in parallel (sixteen H100s at once), each a different
general/domain mix, judged on the \emph{downstream} task. More domain text lowers domain loss but raises
general loss --- a clean tradeoff with a sweet spot at 45--60\% domain.}

\noindent Then the big question: how much domain text? We trained four full models at once, varying only the
mix (30, 45, 60, 75\% domain), and judged them on real pharmaceutical question-answering.

\section{The most important finding in the book}
\fig{figures/fig_overfitting.png}{0.55}{The 75\%-domain model memorized the \emph{style} of medical papers
without learning to answer. Lowest loss, worst answers.}

\noindent The model trained on \textbf{75\% domain} achieved the \emph{best} training loss --- and the
\emph{worst} downstream accuracy. It had over-repeated the thin domain corpus and memorized the surface form
of PubMed abstracts without learning to reason. The balanced models won.

\begin{lesson}[The lesson, in five words]
\textbf{Lower perplexity is not better.} A model's training loss measures how well it predicts text, not how
useful its answers are. Always rank candidates on the real downstream task. We would have shipped the wrong
model if we had trusted loss alone.
\end{lesson}

\section{Why parallel-and-abandon matters}
\fig{figures/fig_ablation.png}{0.5}{Keep the experiments that learn; let the rest fade. Self-abandonment turns
a week of sequential tuning into an afternoon.}

\noindent Running eleven models to choose a recipe sounds expensive. It was not --- the ablation round was
\$15 and the fleet a few hundred dollars --- because the runs are short and the bad ones die early. Tuning one
model at a time would have been slower \emph{and} worse.

\chapter{Instruction tuning (SFT)}
\chapfig{fig_sft.png}

A freshly pretrained model only \emph{continues text}. Ask it a question and it rambles, or asks you a
question back. Supervised fine-tuning (SFT) is the cheap, dramatic step that turns it into something that
\emph{answers}.

\section{How it works}
We format each example with a chat template and --- the single most important detail --- \textbf{mask the loss
to the response tokens only}. The model is never trained to predict the user's question, only the answer.

\begin{lstlisting}
# template:  <|user|>\n{question}<|assistant|>\n{answer}<|endoftext|>
# labels:    [ -1 on the prompt ............ ][ real targets on answer ]
\end{lstlisting}

\noindent Our SFT set was 70{,}000 examples --- medical exam questions with explanations, PubMed-style QA, and
a slice of general instructions so the model stays a generalist.

\section{The before/after}
This is the clearest ``aha'' in the whole project. The same model, on \emph{``What are the side effects of
ibuprofen?''}:

\begin{quote}\small
\textbf{Before SFT (base):} continues into a rambling paragraph that drifts off topic.\\[3pt]
\textbf{After SFT:} ``The common side effects of ibuprofen include: \textbullet\ Headache \textbullet\ Joint
pain \textbullet\ Fever \ldots'' --- a clean, formatted clinical answer.
\end{quote}

\noindent Across our models SFT improved \emph{every} metric, and lifted PubMedQA the most (e.g.\ 0.53
$\rightarrow$ 0.61 on the first model).

\begin{lesson}
Pretraining gives the model \emph{knowledge}; SFT gives it \emph{behaviour}. It is cheap (tens of dollars),
fast, and the most visible single transformation in the pipeline. Mask the loss to the response.
\end{lesson}

\chapter{Evaluation, and the capacity wall}
\chapfig{fig_eval.png}

You cannot improve what you cannot measure. We evaluate with \textbf{answer-likelihood}: for a multiple-choice
question, score each option by the model's length-normalized probability and pick the most likely. This needs
no extra training, so we can track real knowledge throughout.

\section{The numbers, and what they mean}
On the 350M model after SFT: MedMCQA $\approx$0.29 (chance is 0.25), PubMedQA $\approx$0.61. The MedMCQA result
is the one to sit with.

\fig{figures/fig_capacity.png}{0.5}{A 350M model is simply too small to do the multi-step reasoning hard exam
questions demand. More data and more training did not move the score --- it is bounded by capacity.}

\section{``Capacity-bound,'' in plain English}
We \emph{proved} this is a size limit, not a training limit: training five times longer on three times more
data did \textbf{not} raise the MedMCQA score above chance. The bottleneck is the model's size --- its
``brain,'' not its ``study time.'' Yet the same model writes genuinely good, mostly-accurate pharmaceutical
prose, because fluent generation needs far less reasoning capacity than acing trick multiple-choice questions.

\begin{numbers}[The honest scorecard at 350M]
Fluent domain writer: \textbf{yes}. Reliable on open factual questions: \textbf{mostly}. Passes hard medical
exams: \textbf{no} --- that needs scale (Chapter~10) or retrieval (Chapter~11).
\end{numbers}

\begin{lesson}
Match the model size to the \emph{task}. Generation and understanding are cheap; multi-step reasoning is
expensive and scales with parameters. Pick benchmarks and capstones that play to a small model's strengths, or
you will mistake a capacity limit for a failure.
\end{lesson}
